\documentclass{article}
\title{Practice Problem 2.16}
\author{CSAPP}
\usepackage{booktabs}
\usepackage{tabularx}
\usepackage{parskip}
\usepackage{multirow}
\usepackage{array}
\usepackage[table]{xcolor}
\usepackage[margin=1in]{geometry}

\definecolor{headergreen}{RGB}{220, 235, 220}

\begin{document}
	\maketitle

	\begin{flushleft}
		Fill in the table below showing the effects of the different
		shift operations on the single-byte quantities. The best way to think
		about shift operations is to work with binary representation. Convert
		the initial values back to hexadecimal. Each of the answers should
		be 8 binary digits or 2 hexadecimal digits.
	\end{flushleft}

	\begin{center}
		\setlength{\tabcolsep}{4pt}
		\renewcommand{\arraystretch}{2}
		\begin{tabular}{|p{1.0cm}|p{2.1cm}|p{2.1cm}|p{1.0cm}|p{2.1cm}|p{1.0cm}|p{2.1cm}|p{1.0cm}|}
			\hline
			\rowcolor{headergreen}
			\multicolumn{2}{|c|}{$a$}                 &
			\multicolumn{2}{c|}{$a \ll 2$}            &
			\multicolumn{2}{c|}{Logical $a \gg 3$}    &
			\multicolumn{2}{c|}{Arithmetic $a \gg 3$} \\
			\hline
			\texttt{Hex} & \texttt{Binary} & \texttt{Binary} & \texttt{Hex} &
			\texttt{Binary} & \texttt{Hex} & \texttt{Binary} & \texttt{Hex}  \\
			\hline
			\texttt{0xD4} & \texttt{11010100} & \texttt{01010000} & \texttt{0x50} & \texttt{00011010} & \texttt{0x1A} & \texttt{11111010} & \texttt{0xFA} \\
			\hline
			\texttt{0x64} & \texttt{01100100} & \texttt{10010000} & \texttt{0x90} & \texttt{00001100} & \texttt{0x0C} & \texttt{00001100} & \texttt{0x0C} \\
			\hline
			\texttt{0x72} & \texttt{01110010} & \texttt{11001000} & \texttt{0xC8} & \texttt{00001110} & \texttt{0x0E} & \texttt{00001110} & \texttt{0x0E} \\
			\hline
			\texttt{0x44} & \texttt{01000100} & \texttt{00010000} & \texttt{0x10} & \texttt{00001000} & \texttt{0x08} & \texttt{00001000} & \texttt{0x08} \\
			\hline
		\end{tabular}
	\end{center}

\end{document}